\NSPage{097}%
\begin{EESegment}{EE-TGT-P097-001}
On the reserved mean patch, Proposition~\ref{prop:5.5} says that the background still follows the
power law \eqref{eq:4.30}. This form turns the equations into separate power moments of bumps with
compact support. Lemma~\ref{lem:8.7} first builds one fixed linear map that works for any target
values. Lemma~\ref{lem:8.8} then finds the terms that remain when those targets are the defects in
the current fields.
\end{EESegment}

\begin{EESegment}{EE-TGT-P097-002}
For \NSi{NS-F-P097-001}{x=r/\sqrt q\in I_m}, normalizing the velocity at
\NSi{NS-F-P097-002}{q} gives the base
\end{EESegment}
\NSFormula{NS-F-P097-003}{display}{8.24}
\begin{equation}
G_q=0,\qquad V_q=a(\eta)x^{-1-2\lambda},\qquad |a(\eta)|\ge a_0>0.
\tag{8.24}\label{eq:8.24}
\end{equation}
\begin{EESegment}{EE-TGT-P097-003}
Here \NSi{NS-F-P097-004}{a(\eta)=2^{1/2+\lambda}c_{\mathrm{patch}}(1+\eta^2)^{-1}}. Equation
\eqref{eq:5.44} gives the exact statement for the whole summed base. The other tangential
coefficients, beyond the leading profile, have their support elsewhere by \eqref{eq:5.18}.
Multiplying their azimuthal potentials by a cutoff in \NSi{NS-F-P097-005}{q} cannot create an axial
component on this patch, because \NSi{NS-F-P097-006}{q} does not depend on the radius. This means
that \eqref{eq:8.24} holds for the summed base used in the balances. The fixed parameter
\NSi{NS-F-P097-007}{\lambda} is positive.
\end{EESegment}

\begin{lemma}\label{lem:8.7}
\begin{EESegment}{EE-TGT-P097-004}
Suppose the fixed slow base satisfies \eqref{eq:8.24} on the reserved interior interval
\NSi{NS-F-P097-008}{I_m}, where \NSi{NS-F-P097-009}{\lambda>0} and
\NSi{NS-F-P097-010}{|a(\eta)|\ge a_0>0}. It is possible to choose three fixed azimuthal bump
profiles and two fixed axial bump profiles, all supported in this interval, with the following
property. At each \NSi{NS-F-P097-011}{(Z,T)}, take any three scalar targets
\NSi{NS-F-P097-012}{(\mathcal P,\mathcal J_\theta,\mathcal J_z)}. There is then exactly one
linear combination \NSi{NS-F-P097-013}{(\Delta v,\gamma_d)} of the physically rescaled profiles
that satisfies
\end{EESegment}
\NSFormula{NS-F-P097-014}{display}{8.25}
\begin{equation}
\begin{gathered}
\int R^2\Delta v\,dR=0,\qquad \int R\gamma_d\,dR=0,\\
\int (2V/R)\Delta v\,dR=-\mathcal P,\\
\int R^2(G\Delta v+V\gamma_d)\,dR=-\mathcal J_\theta,\\
\int (2RG\gamma_d-RV\Delta v)\,dR=-\mathcal J_z.
\end{gathered}
\tag{8.25}\label{eq:8.25}
\end{equation}
\begin{EESegment}{EE-TGT-P097-005}
Normalize the three targets one row at a time as in \eqref{eq:8.17}. The resulting increments do
not depend on either the angular variables or the auxiliary variables, and they depend linearly on
the targets. For every \NSi{NS-F-P097-015}{\alpha}, targets in
\NSi{NS-F-P097-016}{\mathsf S_\alpha} give increments in
\NSi{NS-F-P097-017}{\mathsf M_\alpha} for every fixed coefficient derivative. The constants in
these bounds may depend on the fixed \NSi{NS-F-P097-018}{\lambda}. The azimuthal-potential
construction in \eqref{eq:8.14} gives \NSi{NS-F-P097-019}{\Delta\gamma=\gamma_d} exactly and
\NSi{NS-F-P097-020}{\Delta\beta\in \mathsf M_{\alpha+1}}. It also preserves both integral
constraints exactly. When the targets are the current defects from Equations~\eqref{eq:8.12} and
\eqref{eq:8.15}, the last three equations cancel both the current defects and the displayed terms
that involve the fixed base.
\end{EESegment}
\end{lemma}

\begin{proof}
\begin{EESegment}{EE-TGT-P097-006}
\textbf{Step 1, choose the five profiles.} Normalize each physical row and its target using
\NSi{NS-F-P097-021}{q}. On this patch, \eqref{eq:8.24} separates the rows into an angular block
with the powers \NSi{NS-F-P097-022}{x} given by
\NSi{NS-F-P097-023}{2,-2-2\lambda,-2\lambda}, and an axial block with the powers
\NSi{NS-F-P097-024}{1,1-2\lambda}. Since the fixed
\NSi{NS-F-P097-025}{\lambda} is positive, the powers in each list are all different.
\end{EESegment}

\begin{EESegment}{EE-TGT-P097-007}
Here is a direct way to see that the map is invertible. Choose a nonnegative bump
\NSi{NS-F-P097-026}{\eta_0}, not identically zero, whose support lies near some
\NSi{NS-F-P097-027}{x_0>0} inside the patch. Next choose a small fixed
\NSi{NS-F-P097-028}{d>0} and make the support narrow enough that the scaled copies
\NSi{NS-F-P097-029}{\eta_j(x)=a_j^{-1}\eta_0(x/a_j),\ a_j=e^{jd},\ j=0,1,2} lie in disjoint
subintervals of the patch. Their power moments are
\end{EESegment}
\NSFormula{NS-F-P097-030}{display}{none}
\[
\int x^p\eta_j(x)\,dx=\mu_p e^{jdp},\qquad
\mu_p=\int x^p\eta_0(x)\,dx>0.
\]
\begin{EESegment}{EE-TGT-P097-008}
For the three angular powers, divide each row of the matrix by
\NSi{NS-F-P097-031}{\mu_p}. What remains is the usual Vandermonde matrix in the different
numbers \NSi{NS-F-P097-032}{e^{dp}}, so its determinant is not zero. Using two copies gives the
same result for the axial block. Every factor left over is a nonzero multiple of
\NSi{NS-F-P097-033}{a(\eta)}.
\end{EESegment}

\NSPage{098}%
\begin{EESegment}{EE-TGT-P098-001}
\textbf{Step 2, find the coefficients by inverting the matrices.} Use
\NSi{NS-F-P098-001}{\eta_0,\eta_1,\eta_2} for the azimuthal profiles and
\NSi{NS-F-P098-002}{\eta_0,\eta_1} for the axial profiles. Write
\NSi{NS-F-P098-003}{\mathcal P_q,\mathcal J_{\theta,q},\mathcal J_{z,q}} for the physical
targets after multiplying them by the powers in \eqref{eq:8.17}, with
\NSi{NS-F-P098-004}{q} in place of \NSi{NS-F-P098-005}{Q}. Define the two matrices
\end{EESegment}
\NSFormula{NS-F-P098-006}{display}{none}
\[
\begin{gathered}
A_\theta=
\begin{pmatrix}
\mu_2 & \mu_2e^{2d} & \mu_2e^{4d}\\
2a\mu_{-2-2\lambda} & 2a\mu_{-2-2\lambda}e^{d(-2-2\lambda)}
  & 2a\mu_{-2-2\lambda}e^{2d(-2-2\lambda)}\\
-a\mu_{-2\lambda} & -a\mu_{-2\lambda}e^{-2d\lambda}
  & -a\mu_{-2\lambda}e^{-4d\lambda}
\end{pmatrix},
\\[4pt]
A_z=
\begin{pmatrix}
\mu_1 & \mu_1e^d\\
a\mu_{1-2\lambda} & a\mu_{1-2\lambda}e^{d(1-2\lambda)}
\end{pmatrix},
\qquad a=a(\eta).
\end{gathered}
\]
\begin{EESegment}{EE-TGT-P098-002}
Step 1 shows that both matrices are invertible. Use their inverses to set
\end{EESegment}
\NSFormula{NS-F-P098-007}{display}{none}
\[
\begin{gathered}
(u_0,u_1,u_2)^T=A_\theta^{-1}(0,-\mathcal P_q,-\mathcal J_{z,q})^T,\\
(s_0,s_1)^T=A_z^{-1}(0,-\mathcal J_{\theta,q})^T,\\
\Delta v_{\mathrm{phys}}(r,z,t)=q^{-A}\sum_{j=0}^{2}u_j\eta_j(r/\sqrt q),\\
\gamma_{d,\mathrm{phys}}(r,z,t)=q^{-A}\sum_{j=0}^{1}s_j\eta_j(r/\sqrt q).
\end{gathered}
\]
\begin{EESegment}{EE-TGT-P098-003}
The matrix \NSi{NS-F-P098-008}{A_\theta} makes the angular-momentum integral zero and gives the
targets \NSi{NS-F-P098-009}{-\mathcal P_q,-\mathcal J_{z,q}}. The matrix
\NSi{NS-F-P098-010}{A_z} makes the axial-flux integral zero and gives the target
\NSi{NS-F-P098-011}{-\mathcal J_{\theta,q}}. Because
\NSi{NS-F-P098-012}{G_q=0}, these two blocks cover all five equations. Their inverses give one
physical correction, depending linearly on the three targets. The correction is smooth in the slow
variables whenever the targets are smooth. The inverse can get worse as
\NSi{NS-F-P098-013}{\lambda\downarrow0}, but no bound uniform in that limit is needed here.
\end{EESegment}

\begin{EESegment}{EE-TGT-P098-004}
\textbf{Step 3, check the normalization and bounds, then build the axial field.} Moving from rows
normalized at \NSi{NS-F-P098-014}{q} to a chart based at a fixed
\NSi{NS-F-P098-015}{Q} multiplies each row and its target by the same powers from
\eqref{eq:8.17}, along with bounded smooth powers of \NSi{NS-F-P098-016}{q/Q}. Every fixed-order
derivative of those ratios is bounded. Inverting and differentiating the matrices therefore costs
only fixed constants. The radial weight is also bounded above and below on the fixed interior
support. This proves the map from \NSi{NS-F-P098-017}{\mathsf S_\alpha} to
\NSi{NS-F-P098-018}{\mathsf M_\alpha}. The required axial field is slow, and its weighted
axial-flux integral is zero. Lemma~\ref{lem:8.2} therefore gives a cutoff remainder that vanishes
identically. Its azimuthal potential and radial velocity component remain inside the mean patch,
including the intervals between the finitely many bump supports.
\end{EESegment}
\end{proof}

\begin{EESegment}{EE-TGT-P098-005}
The linear map is now fixed without using the target values. To use it in the nonlinear equations,
the defects must be recomputed from the updated field. The next lemma gives the exact change and the
extra decay supplied by the bounds used in the iteration. The base estimates needed for that gain
hold on the whole support of the azimuthal potential, including the intervals between the bump
supports.
\end{EESegment}

\begin{lemma}\label{lem:8.8}
\begin{EESegment}{EE-TGT-P098-006}
Suppose the slow base satisfies \eqref{eq:8.24}, with the fixed
\NSi{NS-F-P098-019}{\lambda>0} and the lower bound
\NSi{NS-F-P098-020}{|a(\eta)|\ge a_0>0}. Let
\NSi{NS-F-P098-021}{(\beta,v,\gamma)} be a smooth divergence-free mean correction supported in
the active shell, and keep the base and the wave field \NSi{NS-F-P098-022}{w} fixed. Use
Equations~\eqref{eq:8.3}, \eqref{eq:8.12}, and \eqref{eq:8.15} to form
\NSi{NS-F-P098-023}{g_r} and the three defects
\NSi{NS-F-P098-024}{(\mathcal P,\mathcal J_\theta,\mathcal J_z)}. Feed those three defects into
Lemma~\ref{lem:8.7}, and construct the resulting \NSi{NS-F-P098-025}{\gamma_d} through
\eqref{eq:8.14}. The updated mean velocity is
\end{EESegment}
\NSFormula{NS-F-P098-026}{display}{none}
\[
(\beta_{\mathrm{new}},v_{\mathrm{new}},\gamma_{\mathrm{new}})
=(\beta+\Delta\beta,v+\Delta v,\gamma+\Delta\gamma),
\]
\NSPage{099}%
\begin{EESegment}{EE-TGT-P099-001}
Recompute \NSi{NS-F-P099-001}{g_{r,\mathrm{new}}} from \eqref{eq:8.3}, and define
\NSi{NS-F-P099-002}{\mathcal R_g=g_{r,\mathrm{new}}-g_r-(2V/R)\Delta v}. Then
\end{EESegment}
\NSFormula{NS-F-P099-003}{display}{8.26}
\begin{equation}
\begin{aligned}
\mathcal R_g={}&-\mathfrak t_*\Delta\beta
 -(D_r+1/R)\bigl(2b\Delta\beta+2\beta\Delta\beta+(\Delta\beta)^2\bigr)\\
&-D_z\bigl(b\Delta\gamma+G\Delta\beta+\beta\Delta\gamma
  +\gamma\Delta\beta+\Delta\beta\Delta\gamma\bigr)\\
&+(2v\Delta v+(\Delta v)^2)/R+\varepsilon(\Delta_0-R^{-2})\Delta\beta.
\end{aligned}
\tag{8.26}\label{eq:8.26}
\end{equation}
\begin{EESegment}{EE-TGT-P099-002}
The defects recomputed from the updated field are
\end{EESegment}
\NSFormula{NS-F-P099-004}{display}{8.27}
\begin{equation}
\begin{aligned}
\mathcal P_{\mathrm{new}}&=\int\langle\mathcal R_g\rangle_Y\,dR,\\
(\mathcal J_\theta)_{\mathrm{new}}&=
 \int R^2\langle\gamma\Delta v+v\Delta\gamma+\Delta\gamma\Delta v\rangle_Y\,dR,\\
(\mathcal J_z)_{\mathrm{new}}&=
 \int R\langle2\gamma\Delta\gamma+(\Delta\gamma)^2\rangle_Y\,dR
 -\frac12\int R^2\langle\mathcal R_g\rangle_Y\,dR.
\end{aligned}
\tag{8.27}\label{eq:8.27}
\end{equation}
\begin{EESegment}{EE-TGT-P099-003}
Assume in addition that, everywhere on the support of the azimuthal potential for this correction,
every fixed coefficient derivative of \NSi{NS-F-P099-005}{b} is bounded by
\NSi{NS-F-P099-006}{C_I\varepsilon S_*^{B_I}}, while the corresponding derivatives of the slow
fields \NSi{NS-F-P099-007}{V,G} are bounded by
\NSi{NS-F-P099-008}{C_I S_*^{B_I}}. Suppose the current corrections satisfy
\NSi{NS-F-P099-009}{v,\gamma\in \mathsf M_{0.9}} and
\NSi{NS-F-P099-010}{\beta\in \mathsf M_{1.9}}, and suppose the three targets lie in
\NSi{NS-F-P099-011}{\mathsf S_\alpha} for \NSi{NS-F-P099-012}{\alpha\ge0.9}. Then
\end{EESegment}
\NSFormula{NS-F-P099-013}{display}{none}
\[
\mathcal R_g\in \mathsf M_{\alpha+0.9-2\kappa_s},\qquad
(\mathcal P_{\mathrm{new}},(\mathcal J_\theta)_{\mathrm{new}},
 (\mathcal J_z)_{\mathrm{new}})\in \mathsf S_{\alpha+0.9-2\kappa_s}.
\]
\end{lemma}

\begin{proof}
\begin{EESegment}{EE-TGT-P099-004}
Subtract the two radial sources in \eqref{eq:8.3}. The covariance
\NSi{NS-F-P099-014}{W} does not change. Expanding every quadratic product gives
\eqref{eq:8.26}. In the azimuthal flux defect, the terms linear in the base are
\NSi{NS-F-P099-015}{\int R^2(G\Delta v+V\Delta\gamma)\,dR}. In the axial flux defect, they are
\NSi{NS-F-P099-016}{\int(2RG\Delta\gamma-RV\Delta v)\,dR}. The second term here comes from the
radial-source term \NSi{NS-F-P099-017}{-\frac12\int R^2(2V/R)\Delta v\,dR}. Since
\NSi{NS-F-P099-018}{\Delta\gamma=\gamma_d}, the last three rows of \eqref{eq:8.25} cancel the
old defects together with these linear terms. The products left after that cancellation give
\eqref{eq:8.27}.
\end{EESegment}

\begin{EESegment}{EE-TGT-P099-005}
For the estimates, Lemma~\ref{lem:8.7} gives
\NSi{NS-F-P099-019}{\Delta v,\Delta\gamma\in \mathsf M_\alpha} and
\NSi{NS-F-P099-020}{\Delta\beta\in \mathsf M_{\alpha+1}}. None of these three increments depends
on the auxiliary variable. In particular,
\NSi{NS-F-P099-021}{\mathfrak t_*\Delta\beta=-\varepsilon\partial_T\Delta\beta}. The terms in
\NSi{NS-F-P099-022}{\mathcal R_g} lie in the following classes.
\end{EESegment}
\NSFormula{NS-F-P099-023}{display}{none}
\[
\begin{array}{c|c}
\text{term}&\text{class}\\ \hline
-\mathfrak t_*\Delta\beta&\mathsf M_{\alpha+2}\\
(D_r+1/R)(2b\Delta\beta)&\mathsf M_{\alpha+2-\kappa_s}\\
D_z(b\Delta\gamma+G\Delta\beta)&\mathsf M_{\alpha+2}\\
2v\Delta v/R&\mathsf M_{\alpha+0.9}\\
(\Delta v)^2/R&\mathsf M_{2\alpha}\\
\varepsilon(\Delta_0-R^{-2})\Delta\beta&\mathsf M_{\alpha+2-2\kappa_s}.
\end{array}
\]
\begin{EESegment}{EE-TGT-P099-006}
The other transport products lie in
\NSi{NS-F-P099-024}{\mathsf M_{\alpha+2-\kappa_s}}. Because
\NSi{NS-F-P099-025}{\alpha\ge0.9}, every exponent in the table is at least
\NSi{NS-F-P099-026}{\alpha+0.9-2\kappa_s}. Radial integration keeps these estimates after each
moment has been normalized by the factors given above. The radial weight has a positive lower bound
on the interior support of the azimuthal potential, so the local base bounds in the statement are
enough.
\end{EESegment}
\end{proof}

\begin{EESegment}{EE-TGT-P099-007}
\subsection{Recomputing the fields and matching them between charts.}\label{sec:8.7}
This subsection gives the order for recomputing the residuals and defects in
Equations~\eqref{eq:8.3}, \eqref{eq:8.12}, and \eqref{eq:8.15} after a mean correction. It also
checks that the constructions match between charts, using Equations~\eqref{eq:8.4} and
\eqref{eq:8.23} and Lemma~\ref{lem:6.2}. Each correction changes the quadratic fluxes, so the next
source has to be computed from the updated divergence-free velocity. A cutoff remainder whose
auxiliary average is zero can have a nonzero average after multiplication by another term. Its
contributions to the pressure and to the compatibility integrals in Equations~\eqref{eq:8.3} and
\eqref{eq:8.15} therefore have to be kept in the next step.
\end{EESegment}

\NSPage{100}%
\begin{EESegment}{EE-TGT-P100-001}
Fix the base coefficients and the stress. An actual velocity tuple
\NSi{NS-F-P100-001}{(\beta,v,\gamma,w)} then determines the pressure and the mean residuals in
this order.
\end{EESegment}
\begin{enumerate}[label=(\arabic*)]
\item \begin{EESegment}{EE-TGT-P100-002}
Using the full wave velocities, form
\NSi{NS-F-P100-002}{W_{ab}=\langle w_aw_b\rangle_\theta}. Then define
\NSi{NS-F-P100-003}{g_r} by the third row of \eqref{eq:8.3}.
\end{EESegment}
\item \begin{EESegment}{EE-TGT-P100-003}
Define \NSi{NS-F-P100-004}{\mathcal P=\int\langle g_r\rangle_Y\,dR} and
\NSi{NS-F-P100-005}{p_m=\mathsf T_0(g_r-\rho\mathcal P)}.
\end{EESegment}
\item \begin{EESegment}{EE-TGT-P100-004}
With this pressure, define \NSi{NS-F-P100-006}{E_\theta,E_z} by the first two rows of
\eqref{eq:8.3}.
\end{EESegment}
\item \begin{EESegment}{EE-TGT-P100-005}
Using the same velocity and radial source, define
\NSi{NS-F-P100-007}{\mathcal J_\theta,\mathcal J_z} by \eqref{eq:8.15}.
\end{EESegment}
\end{enumerate}
\begin{EESegment}{EE-TGT-P100-006}
Repeat these definitions after every change to the velocity. For either kind of mean update, the new
tuple is
\end{EESegment}
\NSFormula{NS-F-P100-008}{display}{none}
\[
(\beta,v,\gamma,w)_{\mathrm{new}}
=(\beta+\Delta\beta,v+\Delta v,\gamma+\Delta\gamma,w).
\]
\begin{EESegment}{EE-TGT-P100-007}
For the temporal update, use the current \NSi{NS-F-P100-009}{E_\theta^\circ,E_z^\circ} in
\eqref{eq:8.20}. For the moment update, use the current
\NSi{NS-F-P100-010}{(\mathcal P,\mathcal J_\theta,\mathcal J_z)} in \eqref{eq:8.25}. In both
updates, \NSi{NS-F-P100-011}{\Delta\gamma} means the actual axial increment from
\eqref{eq:8.14}. Its cutoff remainder is therefore included before any of the products in the four
definitions above are evaluated. After the moment update, Equations~\eqref{eq:8.26}--\eqref{eq:8.27}
give the recomputed defects exactly.
\end{EESegment}

\begin{EESegment}{EE-TGT-P100-008}
Proposition~\ref{prop:9.6} and its proof give the correction cycle and the estimates for its
exponents.
\end{EESegment}

\begin{EESegment}{EE-TGT-P100-009}
Because the definitions use physical variables, they match across charts. The radial integrals keep
\NSi{NS-F-P100-012}{(z,t)} fixed, and \NSi{NS-F-P100-013}{q=q(z,t)} keeps them from creating
new dyadic bands. It is therefore possible to fix one common index on a neighborhood that contains
the closures of all the supports in question. The Fourier covariance in \eqref{eq:8.23} makes the
formulas agree where the charts overlap, and the radial formulas agree because they have the same
physical definition. Changing representations shifts the covering indices by only bounded amounts.
The integer-valued covering indices are held fixed while derivatives are taken, and the coordinate
changes preserve the derivative bounds already proved. All the estimates above hold on the same
domain \NSi{NS-F-P100-014}{0<q<q_*}, which was fixed before the iteration. On any closed region
\NSi{NS-F-P100-015}{q\ge c>0}, including \NSi{NS-F-P100-016}{\eta=\pm1}, only finitely many
bands occur. On this region, the fixed-shift integrals, the inverse Fourier transform, and the
finite-dimensional moment correction preserve every one-sided derivative of the source, including
the derivatives at the endpoints.
\end{EESegment}

\begin{EESegment}{EE-TGT-P100-010}
\section{Improving the residual and building the local field}\label{sec:9}
The wave and mean constructions from Sections~\ref{sec:7} and \ref{sec:8} correct different parts
of the momentum residual. This section puts them together to prove Theorem~\ref{thm:3.1}. Start
with the fixed background \NSi{NS-F-P100-017}{(u_B,p_B)} from Proposition~\ref{prop:5.5} and the
primary fields assembled in \eqref{eq:7.35}, which Lemma~\ref{lem:7.7} makes divergence-free. The
construction must produce a local velocity \NSi{NS-F-P100-018}{u_{\mathrm{loc}}} whose residual and every
derivative of that residual vanish to every order at the singularity, while the inner velocity growth
and the exterior heat flow remain unchanged. Every correction creates new nonlinear errors. A full
cycle therefore has to improve the decay before the next cycle begins.
\end{EESegment}

\begin{EESegment}{EE-TGT-P100-011}
One cycle works in the following order. Proposition~\ref{prop:7.2} first removes the nonzero angular
harmonics. Equation~\eqref{eq:7.31} then changes the primary amplitudes to adjust the averaged
stress. The mean equations use \eqref{eq:8.20} to remove the part whose auxiliary average is zero.
Finally, the five-equation system \eqref{eq:8.25} corrects the three compactly supported defects
while preserving the two moment constraints. Proposition~\ref{prop:9.6} proves that a cycle gains a
fixed positive power of \NSi{NS-F-P100-019}{\varepsilon}. Every inverse uses the same fixed
background and the same primary amplitudes. This lets every finite stage live on one common domain,
as Lemma~\ref{lem:9.7} shows. Summing the stages there gives the local field in
Proposition~\ref{prop:9.9}.
\end{EESegment}

\begin{EESegment}{EE-TGT-P100-012}
Table~\ref{tab:1} collects the scales, operators, and field classes that occur repeatedly. In this
section, \NSi{NS-F-P100-020}{\kappa_s=10^{-5}}, and
\NSi{NS-F-P100-021}{\mathsf W_\alpha,\mathsf M_\alpha,\mathsf S_\alpha} have the meaning given in
Section~\ref{sec:6.4} at every derivative order. Each exponent is a power of
\NSi{NS-F-P100-022}{\varepsilon=Q^h}, with the normalization
\end{EESegment}
\NSPage{101}%
\begin{EESegment}{EE-TGT-P101-001}
used for those classes. At a fixed correction stage and a fixed order of amplitude differentiation,
the constants, the powers of \NSi{NS-F-P101-001}{S_*}, and the inverse powers of distance are
allowed to depend on that stage and order. They are uniform in the band, label, copy, and point. The
angular mean \NSi{NS-F-P101-002}{\langle\cdot\rangle_\theta} and the auxiliary mean
\NSi{NS-F-P101-003}{\langle\cdot\rangle_Y} are different operations, both defined in
\eqref{eq:3.7}. Every class estimate is an estimate for normalized chart coefficients. The factors
that recover the physical velocity, pressure, and residual are
\NSi{NS-F-P101-004}{Q^{-A}}, \NSi{NS-F-P101-005}{Q^{-2A}}, and
\NSi{NS-F-P101-006}{Q^{-2A-1/2}}, respectively. In particular, the physical fields have residual
\end{EESegment}
\NSFormula{NS-F-P101-007}{display}{none}
\[
\mathscr R(u,p)=\partial_tu+(u\cdot\nabla)u-\Delta u+\nabla p,
\qquad
\mathscr R_*=Q^{2A+1/2}\mathscr R(u,p).
\]
\begin{EESegment}{EE-TGT-P101-002}
The residual identities below use physical fields, while their estimates use normalized
coefficients. A Cartesian space--time jet collects a field and all its Cartesian space--time
derivatives up to the stated order. Equation~\eqref{eq:5.23} writes
\NSi{NS-F-P101-008}{|\cdot|_m} for the largest absolute value among them through that order.
\end{EESegment}

\begin{EESegment}{EE-TGT-P101-003}
\subsection{Residual estimates for velocities built from curls}\label{sec:9.1}
The amplitude equations \eqref{eq:7.5} cancel the main part of a prescribed source. The physical
velocity also contains the terms created by taking a curl, as Lemma~\ref{lem:7.7} describes. Its
residual contains derivatives of the amplitudes and the background as well. Equation~\eqref{eq:9.2}
expands those terms. Proposition~\ref{prop:9.1} and Lemma~\ref{lem:9.2} first estimate the linear
terms left after the main cancellation and the nonlinear interactions. Those estimates give the
gains needed when Proposition~\ref{prop:9.6} adds one correction after another.
\end{EESegment}

\begin{EESegment}{EE-TGT-P101-004}
For a fixed label, write an angular harmonic as
\NSi{NS-F-P101-009}{a\exp(ikm\Phi)}, where
\NSi{NS-F-P101-010}{m\in\mathbb Z\setminus\{0\}} and
\NSi{NS-F-P101-011}{a} is the amplitude. Equation~\eqref{eq:7.4} defines the phase normal
\NSi{NS-F-P101-012}{n_\Phi}, and \eqref{eq:6.6} gives the chain-rule operators for evaluation at
\NSi{NS-F-P101-013}{Y=Y(r,t)}. The background coefficients
\NSi{NS-F-P101-014}{(b,V,G)} are the ones used in Sections~\ref{sec:6} and \ref{sec:7}.
\end{EESegment}

\begin{EESegment}{EE-TGT-P101-005}
For a transverse amplitude \NSi{NS-F-P101-015}{t_m}, the potential in \eqref{eq:7.38} gives
\end{EESegment}
\NSFormula{NS-F-P101-016}{display}{none}
\[
C_m=\frac{i n_\Phi\times t_m}{km|n_\Phi|^2},
\qquad
\operatorname{curl}_*(C_me^{ikm\Phi})=(t_m+r_m)e^{ikm\Phi},
\qquad
a_m=t_m+r_m.
\]
\begin{EESegment}{EE-TGT-P101-006}
Here \NSi{NS-F-P101-017}{r_m} is the correction created by taking the curl. If
\NSi{NS-F-P101-018}{t_m\in \mathsf W_\alpha}, then Lemma~\ref{lem:7.7} gives
\NSi{NS-F-P101-019}{r_m\in \mathsf W_{\alpha+1/2-\kappa_s}}. Here longitudinal means measured in
the direction of the phase normal. Equation~\eqref{eq:7.39} estimates the scalar product with that normal for the
full amplitude \NSi{NS-F-P101-020}{a_m}. This difference matters
in the quadratic transport term, because the factor doing the advecting is the full divergence-free
velocity.
\end{EESegment}

\begin{EESegment}{EE-TGT-P101-007}
Separate the principal linear operator, which the amplitude equation handles, from all the other
transport, pressure, and viscous terms. For comparison with the residual later, write the principal
operator acting on an amplitude and its pressure as
\end{EESegment}
\NSFormula{NS-F-P101-021}{display}{9.1}
\begin{equation}
\mathscr L_m(a,\pi)=Q^{1+h}N_{\mathrm{abs}}a+\mathcal Ka
+\varepsilon k^2m^2|n_\Phi|^2a+ikmn_\Phi\pi.
\tag{9.1}\label{eq:9.1}
\end{equation}
\begin{EESegment}{EE-TGT-P101-008}
The inverse from Proposition~\ref{prop:7.2} constructs
\NSi{NS-F-P101-022}{(t_m,\pi_m)} so that
\NSi{NS-F-P101-023}{\mathscr L_m(t_m,\pi_m)=-f_m} and
\NSi{NS-F-P101-024}{n_\Phi\cdot t_m=0}, before the pulse cutoff is applied. The amplitude
\NSi{NS-F-P101-025}{t_m} does not have to extend by zero in time before that cutoff is applied.
\end{EESegment}

\begin{proposition}\label{prop:9.1}
\begin{EESegment}{EE-TGT-P101-009}
Fix a label, \NSi{NS-F-P101-026}{m\in\mathbb Z\setminus\{0\}}, and
\NSi{NS-F-P101-027}{\alpha\in\mathbb R}. Let \NSi{NS-F-P101-028}{f_m\in \mathsf W_\alpha}
satisfy the support and smooth-extension assumptions in Proposition~\ref{prop:7.2}. Also assume
that the source harmonic \NSi{NS-F-P101-029}{f_me^{ikm\Phi}} extends smoothly by zero outside the
fixed local rectangle, including across the ends of its pulse. Let
\NSi{NS-F-P101-030}{(t_m,\pi_m)} solve
\NSi{NS-F-P101-031}{\mathscr L_m(t_m,\pi_m)=-f_m}. The same conclusion holds for a homogeneous
pulse \NSi{NS-F-P101-032}{t_m\in \mathsf W_\alpha}, with
\NSi{NS-F-P101-033}{f_m=0}. Multiply the potential and the pressure by the pulse cutoff, then take
the curl as in \eqref{eq:7.38}. After the prescribed source
\NSi{NS-F-P101-034}{f_me^{ikm\Phi}} has been added, the normalized linear residual lies in
\NSi{NS-F-P101-035}{\mathsf W_{\alpha+1/2-3\kappa_s}}. The only terms left outside that class are
additive tails from the pulse cutoff, and every Cartesian space--time derivative of each tail
vanishes to infinite order as \NSi{NS-F-P101-036}{q\downarrow0}. The pressure amplitude
\NSi{NS-F-P101-037}{\pi_m} lies in \NSi{NS-F-P101-038}{\mathsf W_{\alpha+1/2}}.
\end{EESegment}
\end{proposition}

\begin{proof}
\begin{EESegment}{EE-TGT-P101-010}
Expand the linearized residual into its principal operator \eqref{eq:9.1} and the terms left over.
For this calculation, \NSi{NS-F-P101-039}{a} is any harmonic amplitude and
\NSi{NS-F-P101-040}{\pi} is its pressure amplitude. Let
\end{EESegment}
\NSPage{102}%
\begin{EESegment}{EE-TGT-P102-001}
\NSi{NS-F-P102-001}{E_{ik}=(\mathfrak t_*+bD_r+(V/R)\partial_\theta+GD_z)\Phi} be the phase
transport defect from \eqref{eq:7.9}. The terms left after linearizing around
\NSi{NS-F-P102-002}{(b,V,G)} are
\end{EESegment}
\NSFormula{NS-F-P102-003}{display}{9.2}
\begin{equation}
\begin{aligned}
\mathscr L_m^{\mathrm{rem}}(a,\pi)={}&
 (-\varepsilon\partial_T+bD_r+GD_z)a+ikmE_{ik}a\\
&+a_r(\partial_Rb)e_r+a_zD_z(b,V,G)+(b/R)a_\theta e_\theta
 +(D_r\pi,0,D_z\pi)\\
&-\varepsilon\left\{
\begin{aligned}
 &(D_r^2+R^{-1}D_r+D_z^2)a-R^{-2}(a_r,a_\theta,0)\\
 &\quad+2ikm(n_{\Phi,r}D_r+n_{\Phi,z}D_z)a\\
 &\quad+ikm(D_rn_{\Phi,r}+n_{\Phi,r}/R+D_zn_{\Phi,z})a\\
 &\quad+2ikm(n_{\Phi,\theta}/R)(-a_\theta,a_r,0)
\end{aligned}
\right\}.
\end{aligned}
\tag{9.2}\label{eq:9.2}
\end{equation}
\begin{EESegment}{EE-TGT-P102-002}
For a cylindrical vector \NSi{NS-F-P102-004}{v} transported by another cylindrical vector
\NSi{NS-F-P102-005}{u}, differentiation of the moving cylindrical basis contributes
\NSi{NS-F-P102-006}{(-u_\theta v_\theta,u_\theta v_r,0)/R}. Linearizing that term gives the
connections in \NSi{NS-F-P102-007}{\mathcal K} and the displayed
\NSi{NS-F-P102-008}{b/R} term. Applying the scalar Laplacian to the harmonic gives the quadratic
phase-gradient term in \eqref{eq:9.1} and the two mixed phase--amplitude derivative terms in
\eqref{eq:9.2}. Applying the vector Laplacian also gives the last angular connection and the
\NSi{NS-F-P102-009}{R^{-2}} term. The amplitude coefficients do not depend on
\NSi{NS-F-P102-010}{\theta}.
\end{EESegment}

\begin{EESegment}{EE-TGT-P102-003}
The transverse-amplitude equation in Proposition~\ref{prop:7.2} gives
\NSi{NS-F-P102-011}{\mathscr L_m(t_m,\pi_m)=-f_m}, with
\NSi{NS-F-P102-012}{\pi_m\in \mathsf W_{\alpha+1/2}}. For each kind of term in
\eqref{eq:9.2}, the gain over \NSi{NS-F-P102-013}{\alpha} is listed below.
\end{EESegment}
\NSFormula{NS-F-P102-014}{display}{none}
\[
\begin{array}{c|c}
\text{term}&\text{gain}\\ \hline
\text{slow transport}&1-\kappa_s\\
\text{phase transport defect}&1/2\\
\text{base derivatives and connections}&1\\
\text{pressure-amplitude gradient}&1/2-\kappa_s\\
\text{viscous amplitude derivatives}&1-2\kappa_s\\
\text{mixed derivatives and phase divergence}&1/2-\kappa_s\\
\text{viscous angular connection}&1/2
\end{array}
\]
\begin{EESegment}{EE-TGT-P102-004}
The bound for the phase defect allows polynomial factors of
\NSi{NS-F-P102-015}{S_*}. Applying the principal velocity operator to the curl correction does not
change its exponent. The fast-time derivative differentiates an amplitude coefficient, and
\NSi{NS-F-P102-016}{\varepsilon k^2=O(1)}. At every derivative order, the potential formula and
these chain-rule operators give the common lower bound
\NSi{NS-F-P102-017}{\alpha+1/2-3\kappa_s} for the exponents of these terms.
\end{EESegment}

\begin{EESegment}{EE-TGT-P102-005}
Apply the pulse cutoff to the potential and the pressure before taking the curl. The source left
uncancelled, \NSi{NS-F-P102-018}{(1-\psi)f}, and the main derivative of the pulse cutoff,
\NSi{NS-F-P102-019}{\psi't}, together with its fixed multipliers, are supported where
\NSi{NS-F-P102-020}{P_v\le e^{-cS_*}}. For every fixed physical derivative, their bounds have the
form \NSi{NS-F-P102-021}{CQ^{-M}S_*^Pe^{-cS_*}}. They are therefore flat, because
\NSi{NS-F-P102-022}{S_*=\ell^2} and \NSi{NS-F-P102-023}{Q=2^{-\ell}}. Keep these terms in the
additive residual throughout the iteration. Later forward solves are applied only to the supported
source terms.
\end{EESegment}
\end{proof}

\begin{EESegment}{EE-TGT-P102-006}
The nonlinear terms need separate estimates for interactions between a wave and a mean field and
for interactions between two waves. In the two-wave case, incompressibility controls the contraction
of the advecting amplitude with the phase gradient.
\end{EESegment}

\begin{EESegment}{EE-TGT-P102-007}
In the next lemma, \NSi{NS-F-P102-024}{(w\cdot\nabla_*)v} means cylindrical transport with the
normalized derivatives \NSi{NS-F-P102-025}{(D_r,R^{-1}\partial_\theta,D_z)}, including the terms
created by differentiating the cylindrical frame. The lemma therefore estimates the normalized
nonlinear residual.
\end{EESegment}

\begin{lemma}\label{lem:9.2}
\begin{EESegment}{EE-TGT-P102-008}
Let \NSi{NS-F-P102-026}{w} be a wave in \NSi{NS-F-P102-027}{\mathsf W_\alpha}. Let
\NSi{NS-F-P102-028}{v} be a mean field whose tangential components lie in
\NSi{NS-F-P102-029}{\mathsf M_\mu} and whose radial component lies in
\NSi{NS-F-P102-030}{\mathsf M_{\mu+1}}. Then the nonzero harmonics of
\NSi{NS-F-P102-031}{(w\cdot\nabla_*)v+(v\cdot\nabla_*)w} lie in
\NSi{NS-F-P102-032}{\mathsf W_{\alpha+\mu-1/2}}.
\end{EESegment}
